%%%%%%%%%%%%%%%%%%%%%%%%%%%%%%%%%%%%%%%%%%%%%%%%%%%%%%%%%%%%
%% Chapter 1 — Introduction (indicative length: 6–9 pages)
%% Purpose: introduce the practical and scientific problem,
%% motivate the proposed learning-based adaptive architecture,
%% and state exactly what the thesis contributes.
%% Main output: research problem, research questions,
%% objectives, contributions, scope, and chapter roadmap.
%%%%%%%%%%%%%%%%%%%%%%%%%%%%%%%%%%%%%%%%%%%%%%%%%%%%%%%%%%%%
\chapter{Introduction}
\label{chap:introduction}

\section{Motivation and Application Context}
\label{sec:intro:motivation}
% - Heavy-duty and mobile manipulators; why hydraulic actuation is used
%   when high force density and robustness are required.
% - Difficulty of precise end-effector tracking with nonlinear,
%   load-dependent, uncertain actuator response not captured by a
%   purely kinematic controller.
% - Motivate a structured approach: keep a nominal task-space
%   controller, learn the actuator behaviour and residual corrections.
\todo{write motivation}

\section{Problem Statement}
\label{sec:intro:problem}
% - Gap between task-space motion commands and the valve/cylinder
%   behaviour required to realize them.
% - Why a fixed nominal model fails when hydraulic or mechanical
%   parameters vary.
% - Need for a model compact enough for RL simulation but expressive
%   enough to represent actuator variability.
\todo{write problem statement}

\section{Research Questions}
\label{sec:intro:research-questions}
% RQ1: Can learned cylinder-space dynamics improve task-space tracking
%      combined with a nominal DLS controller and residual PPO?
% RQ2: Does training on randomized Simscape data improve robustness
%      compared with nominal-only model training?
% RQ3: Does an AAM-estimated latent representation improve the ROM
%      beyond robust training without online adaptation?
% RQ4: How does the proposed ROM-AAM pipeline compare with the learned
%      ID-FD pipeline under matched control and test conditions?
\todo{state research questions RQ1--RQ4}

\section{Objectives}
\label{sec:intro:objectives}
% - Generate informative nominal and randomized datasets from a
%   high-fidelity Simscape model.
% - Train and validate cylinder-space ID, FD, and reduced-order models.
% - Learn a latent representation of randomized system dynamics and
%   estimate it online with an AAM.
% - Integrate PPO1 and PPO2 pipelines into IsaacLab and evaluate
%   tracking, robustness, adaptation, and computational behaviour.
\todo{state objectives}

\section{Contributions}
\label{sec:intro:contributions}
% - Combined closed-loop and open-loop data-generation strategy for
%   learning cylinder-space actuator dynamics.
% - Learned ID-FD control pipeline as a non-adaptive learning-based
%   reference architecture.
% - Latent-conditioned reduced-order model and AAM-based online
%   adaptation mechanism.
% - Residual PPO task-space control framework integrating classical
%   kinematics with learned actuator models.
% - Controlled comparison across nominal, randomized/robust, and
%   adaptive configurations.
\todo{list contributions}

\section{Scope and Limitations}
\label{sec:intro:scope}
% - Simulation-only vs. real-system validation.
% - Which arm DOFs and cylinders are modeled (incl. passive coupling
%   of the fifth joint).
% - Latent representation is not automatically equivalent to
%   identifiable physical parameters.
% - Safety, hardware, and real-time deployment questions outside scope.
\todo{define scope and limitations}

\section{Thesis Structure}
\label{sec:intro:structure}
% One concise paragraph per chapter (Chapters 2--8).
\todo{summarize chapters 2--8}
